% --- Archon physics formalization source begin ---
% archon:physics
% archon:covers PhyXMiniProblems/problem_phyx_mini_0683.lean
% archon:source-report reports/phyx_mini/problem_phyx_mini_0683.source.json

\chapter{Physics problem phyx\_mini\_0683}
\label{ch:PhyXMiniProblems_problem_phyx_mini_0683}

\paragraph{Problem source.}
\#\# Physical scenario

A firefighter, a distance \textbackslash{}( d \textbackslash{}) from a burning building, directs a stream of water from a fire hose at angle \textbackslash{}( \textbackslash{}theta\_i \textbackslash{}) above the horizontal as shown in figure.

\#\# Question

If the initial speed of the stream is \textbackslash{}( v\_i \textbackslash{}), at what height \textbackslash{}( h \textbackslash{}) does the water strike the building?

\#\# Answer choices

- A: A: \textbackslash{}( d \textbackslash{}tan \textbackslash{}theta\_i - \textbackslash{}frac\{gd\textasciicircum{}2\}\{2v\_i\textasciicircum{}2 \textbackslash{}cot\textasciicircum{}2 \textbackslash{}theta\_i\} \textbackslash{})

- B: B: \textbackslash{}( d \textbackslash{}tan \textbackslash{}theta\_i - \textbackslash{}frac\{gd\textasciicircum{}2\}\{2v\_i\textasciicircum{}2 \textbackslash{}tan\textasciicircum{}2 \textbackslash{}theta\_i\} \textbackslash{})

- C: C: \textbackslash{}( d \textbackslash{}tan \textbackslash{}theta\_i - \textbackslash{}frac\{gd\textasciicircum{}2\}\{2v\_i\textasciicircum{}2 \textbackslash{}cos\textasciicircum{}2 \textbackslash{}theta\_i\} \textbackslash{})

- D: D: \textbackslash{}( d \textbackslash{}tan \textbackslash{}theta\_i - \textbackslash{}frac\{gd\textasciicircum{}2\}\{2v\_i\textasciicircum{}2 \textbackslash{}sin\textasciicircum{}2 \textbackslash{}theta\_i\} \textbackslash{})

\#\# Figure caption (auxiliary; use the image as primary evidence)

The image depicts a firefighting scene with the following elements:

1. **Firefighter**: A person in orange firefighting gear is holding a hose, directing a water stream.

2. **Water Stream**: The water, depicted as a blue, arcing line, is being projected from the hose towards a burning building, following a parabolic trajectory.

3. **Building on Fire**: A multi-story building to the right is emitting flames and smoke, which are shown coming from a window.

4. **Vectors and Angles**:
   - A red vector labeled \textbackslash{}(\textbackslash{}vec\{v\}\_i\textbackslash{}) represents the initial velocity of the water stream.
   - An angle \textbackslash{}(\textbackslash{}theta\_i\textbackslash{}) is shown between the ground and the initial velocity vector, indicating the launch angle of the water stream.

5. **Distances**:
   - The horizontal distance from the firefighter to the building is labeled as \textbackslash{}(d\textbackslash{}).
   - The vertical distance from the ground to the point where the water hits the building is labeled as \textbackslash{}(h\textbackslash{}).

This illustration likely serves as an example of projectile motion, demonstrating the application of physics in firefighting scenarios.

\#\# Dataset metadata

Kinematics; Physical Model Grounding Reasoning, Spatial Relation Reasoning

\paragraph{Recorded answer/context.}
C: C: \textbackslash{}( d \textbackslash{}tan \textbackslash{}theta\_i - \textbackslash{}frac\{gd\textasciicircum{}2\}\{2v\_i\textasciicircum{}2 \textbackslash{}cos\textasciicircum{}2 \textbackslash{}theta\_i\} \textbackslash{})

\paragraph{Figure/image path.}
/root/proposal\_for\_physic/science-mango/phyx\_mini\_run/phyx\_data/test\_image/683.png

\paragraph{Formalization target.}
create a compiling Lean file with sorry bodies at `PhyXMiniProblems/problem\_phyx\_mini\_0683.lean`. The Lean declarations must preserve the physical quantities, dimensions or dimensional roles, figure labels, governing-law hypotheses, and final relation expressed by this problem.
Use Mathlib/Physlib names found through LeanExplore where available. If a physics API is missing, introduce faithful local abstractions rather than scalar placeholder aliases.

\begin{theorem}[Physics formalization target]
\label{thm:physics:phyx_mini_0683:target}
\lean{PhyXMiniProblems.ProblemPhyXMini0683.problem_phyx_mini_0683}
\uses{def:physics:phyx-mini-0683:phyxminiproblems-problemphyxmini0683-lengthmagnitudeinmeters, def:physics:phyx-mini-0683:phyxminiproblems-problemphyxmini0683-speedinmeterspersecond, def:physics:phyx-mini-0683:phyxminiproblems-problemphyxmini0683-accelerationinmeterspersecondsquared, def:physics:phyx-mini-0683:phyxminiproblems-problemphyxmini0683-firehoseprojectilesetup, def:physics:phyx-mini-0683:phyxminiproblems-problemphyxmini0683-satisfiesuniformgravityprojectilemotion, def:physics:phyx-mini-0683:phyxminiproblems-problemphyxmini0683-matchesfirehosefigure}
The assigned autoformalize agent should translate this physics problem into Lean declarations in the covered file, with theorem and lemma proof bodies written as `by sorry`.
\end{theorem}
\begin{proof}
This is an autoformalization task, not a proof task. Produce statements that can later be proved by the physics prover without weakening the physical model.
\end{proof}

% --- Archon declaration topology begin ---
\section{Declaration topology}
\begin{definition}[acceleration Dimension]
  \label{def:physics:phyx-mini-0683:phyxminiproblems-problemphyxmini0683-accelerationdimension}
  \lean{PhyXMiniProblems.ProblemPhyXMini0683.accelerationDimension}
  Acceleration has physical dimension length divided by time squared
\end{definition}
\begin{proof}
  This is the declared model definition of acceleration Dimension.
\end{proof}
\begin{definition}[Length Magnitude]
  \label{def:physics:phyx-mini-0683:phyxminiproblems-problemphyxmini0683-lengthmagnitude}
  \lean{PhyXMiniProblems.ProblemPhyXMini0683.LengthMagnitude}
  A nonnegative physical length, used for the figure labels d and h
\end{definition}
\begin{proof}
  This is the declared model definition of Length Magnitude.
\end{proof}
\begin{definition}[Signed Length Quantity]
  \label{def:physics:phyx-mini-0683:phyxminiproblems-problemphyxmini0683-signedlengthquantity}
  \lean{PhyXMiniProblems.ProblemPhyXMini0683.SignedLengthQuantity}
  A signed physical length, used for Cartesian coordinate components
\end{definition}
\begin{proof}
  This is the declared model definition of Signed Length Quantity.
\end{proof}
\begin{definition}[Acceleration Magnitude]
  \label{def:physics:phyx-mini-0683:phyxminiproblems-problemphyxmini0683-accelerationmagnitude}
  \lean{PhyXMiniProblems.ProblemPhyXMini0683.AccelerationMagnitude}
  \uses{def:physics:phyx-mini-0683:phyxminiproblems-problemphyxmini0683-accelerationdimension}
  A nonnegative physical acceleration magnitude, used for g
\end{definition}
\begin{proof}
  This is the declared model definition of Acceleration Magnitude.
\end{proof}
\begin{definition}[length Magnitude In Meters]
  \label{def:physics:phyx-mini-0683:phyxminiproblems-problemphyxmini0683-lengthmagnitudeinmeters}
  \lean{PhyXMiniProblems.ProblemPhyXMini0683.lengthMagnitudeInMeters}
  \uses{def:physics:phyx-mini-0683:phyxminiproblems-problemphyxmini0683-lengthmagnitude}
  Read a nonnegative physical length in metres
\end{definition}
\begin{proof}
  This is the declared model definition of length Magnitude In Meters.
\end{proof}
\begin{definition}[signed Length In Meters]
  \label{def:physics:phyx-mini-0683:phyxminiproblems-problemphyxmini0683-signedlengthinmeters}
  \lean{PhyXMiniProblems.ProblemPhyXMini0683.signedLengthInMeters}
  \uses{def:physics:phyx-mini-0683:phyxminiproblems-problemphyxmini0683-signedlengthquantity}
  Read a signed Cartesian length component in metres
\end{definition}
\begin{proof}
  This is the declared model definition of signed Length In Meters.
\end{proof}
\begin{definition}[speed In Meters Per Second]
  \label{def:physics:phyx-mini-0683:phyxminiproblems-problemphyxmini0683-speedinmeterspersecond}
  \lean{PhyXMiniProblems.ProblemPhyXMini0683.speedInMetersPerSecond}
  Read a nonnegative speed in metres per second
\end{definition}
\begin{proof}
  This is the declared model definition of speed In Meters Per Second.
\end{proof}
\begin{definition}[acceleration In Meters Per Second Squared]
  \label{def:physics:phyx-mini-0683:phyxminiproblems-problemphyxmini0683-accelerationinmeterspersecondsquared}
  \lean{PhyXMiniProblems.ProblemPhyXMini0683.accelerationInMetersPerSecondSquared}
  \uses{def:physics:phyx-mini-0683:phyxminiproblems-problemphyxmini0683-accelerationmagnitude}
  Read a nonnegative acceleration in metres per second squared
\end{definition}
\begin{proof}
  This is the declared model definition of acceleration In Meters Per Second Squared.
\end{proof}
\begin{definition}[Planar Position]
  \label{def:physics:phyx-mini-0683:phyxminiproblems-problemphyxmini0683-planarposition}
  \lean{PhyXMiniProblems.ProblemPhyXMini0683.PlanarPosition}
  \uses{def:physics:phyx-mini-0683:phyxminiproblems-problemphyxmini0683-signedlengthquantity}
  A point in the vertical plane of the illustrated water trajectory
\end{definition}
\begin{proof}
  This is the declared model definition of Planar Position.
\end{proof}
\begin{definition}[Figure Point]
  \label{def:physics:phyx-mini-0683:phyxminiproblems-problemphyxmini0683-figurepoint}
  \lean{PhyXMiniProblems.ProblemPhyXMini0683.FigurePoint}
  The three geometrically significant points marked or implied by the image
\end{definition}
\begin{proof}
  This is the declared model definition of Figure Point.
\end{proof}
\begin{definition}[Fire Hose Figure]
  \label{def:physics:phyx-mini-0683:phyxminiproblems-problemphyxmini0683-firehosefigure}
  \lean{PhyXMiniProblems.ProblemPhyXMini0683.FireHoseFigure}
  \uses{def:physics:phyx-mini-0683:phyxminiproblems-problemphyxmini0683-lengthmagnitude, def:physics:phyx-mini-0683:phyxminiproblems-problemphyxmini0683-planarposition, def:physics:phyx-mini-0683:phyxminiproblems-problemphyxmini0683-figurepoint}
  Geometric objects and labels read from the supplied firefighting figure. The booleans retain visible spatial information that does not enter the final closed form directly
\end{definition}
\begin{proof}
  This is the declared model definition of Fire Hose Figure.
\end{proof}
\begin{definition}[Fire Hose Projectile Setup]
  \label{def:physics:phyx-mini-0683:phyxminiproblems-problemphyxmini0683-firehoseprojectilesetup}
  \lean{PhyXMiniProblems.ProblemPhyXMini0683.FireHoseProjectileSetup}
  \uses{def:physics:phyx-mini-0683:phyxminiproblems-problemphyxmini0683-lengthmagnitude, def:physics:phyx-mini-0683:phyxminiproblems-problemphyxmini0683-accelerationmagnitude, def:physics:phyx-mini-0683:phyxminiproblems-problemphyxmini0683-planarposition, def:physics:phyx-mini-0683:phyxminiproblems-problemphyxmini0683-firehosefigure}
  The independent physical quantities and observables of the fire-hose stream. The trajectory and impact observables are not defined from the requested height formula; they are related only by the governing-law and figure premises below
\end{definition}
\begin{proof}
  This is the declared model definition of Fire Hose Projectile Setup.
\end{proof}
\begin{definition}[Satisfies Uniform Gravity Projectile Motion]
  \label{def:physics:phyx-mini-0683:phyxminiproblems-problemphyxmini0683-satisfiesuniformgravityprojectilemotion}
  \lean{PhyXMiniProblems.ProblemPhyXMini0683.SatisfiesUniformGravityProjectileMotion}
  \uses{def:physics:phyx-mini-0683:phyxminiproblems-problemphyxmini0683-signedlengthinmeters, def:physics:phyx-mini-0683:phyxminiproblems-problemphyxmini0683-speedinmeterspersecond, def:physics:phyx-mini-0683:phyxminiproblems-problemphyxmini0683-accelerationinmeterspersecondsquared, def:physics:phyx-mini-0683:phyxminiproblems-problemphyxmini0683-firehoseprojectilesetup}
  The no-drag constant-gravity projectile equations, expressed in coherent SI component readouts for every nonnegative time in seconds. The launch point is the coordinate origin, so the two initial-position terms are zero
\end{definition}
\begin{proof}
  This is the declared model definition of Satisfies Uniform Gravity Projectile Motion.
\end{proof}
\begin{definition}[Matches Fire Hose Figure]
  \label{def:physics:phyx-mini-0683:phyxminiproblems-problemphyxmini0683-matchesfirehosefigure}
  \lean{PhyXMiniProblems.ProblemPhyXMini0683.MatchesFireHoseFigure}
  \uses{def:physics:phyx-mini-0683:phyxminiproblems-problemphyxmini0683-lengthmagnitudeinmeters, def:physics:phyx-mini-0683:phyxminiproblems-problemphyxmini0683-signedlengthinmeters, def:physics:phyx-mini-0683:phyxminiproblems-problemphyxmini0683-firehoseprojectilesetup}
  Primary-image readouts and incidence geometry. This connects the labels d, h, v\_i, and theta\_i to the physical setup, places the nozzle and building on the same horizontal ground, and identifies the positive-time trajectory point that strikes the building. It contains no eliminated height formula
\end{definition}
\begin{proof}
  This is the declared model definition of Matches Fire Hose Figure.
\end{proof}
% --- Archon declaration topology end ---
% --- Archon physics formalization source end ---
